% =============================================================================
\chapter{Phân tích thiết kế}\label{chap:thietke}
% =============================================================================

\section{Mô hình dữ liệu}

Phân tích bài toán cho bốn thực thể chính:

\begin{itemize}
  \item \textbf{Mặt hàng} (\code{MatHang}): mã vạch (khóa), tên, nhóm, đơn giá bán, ngưỡng cảnh báo,
        tồn kho và danh sách các lô còn hàng.
  \item \textbf{Lô} (\code{Lo}): một lần nhập của một mặt hàng -- mã lô, ngày nhập, \emph{hạn sử dụng},
        số lượng \emph{còn lại}, giá nhập. HSD gắn với lô chứ không gắn với mặt hàng, vì mỗi lần nhập có HSD khác nhau.
  \item \textbf{Phiếu} (\code{Phieu}): một giao dịch nhập hoặc xuất -- mã phiếu, loại, ngày, mã vạch,
        số lượng, đơn giá tại thời điểm giao dịch. Phiếu là lịch sử, không bị xóa khi mặt hàng bị xóa.
  \item \textbf{Kho} (\code{Kho}): gồm chỉ mục mặt hàng (B-Tree), danh sách phiếu và bộ đếm cấp mã lô.
\end{itemize}

Hai quyết định thiết kế quan trọng được cả nhóm thống nhất từ đầu:

\begin{enumerate}
  \item \textbf{Tồn kho là dữ liệu dẫn xuất.} \code{tonKho} luôn bằng tổng \code{soLuong} các lô và
        \emph{không ghi ra tệp}; khi đọc tệp, chương trình cộng số lượng các lô để tính lại. Nhờ vậy tồn kho không
        bao giờ lệch với các lô.
  \item \textbf{Quyền sở hữu bộ nhớ rõ ràng.} B-Tree chỉ giữ con trỏ \code{MatHang*}; \code{btRemove} gỡ mặt hàng
        khỏi cây và trả con trỏ về, còn mô-đun \code{Inventory} là nơi duy nhất \code{delete} mặt hàng cùng các lô.
        Quy ước này tránh giải phóng hai lần.
\end{enumerate}

Mã nguồn~\ref{lst:types} trích các kiểu dữ liệu chính trong \code{Types.h}.

\begin{lstlisting}[float=tbp,caption={Các kiểu dữ liệu chính trong \code{Types.h} (trích)},label=lst:types]
struct Date { int d; int m; int y; };

struct Lo {                       // mot lan nhap cua mot mat hang
    int maLo;  Date ngayNhap;  Date hsd;   // HSD gan voi tung lo
    int soLuong;                  // con lai trong lo
    double giaNhap;  Lo* next;
};
struct LotQueue { Lo* front; Lo* rear; int count; };   // front = lo cu nhat

struct MatHang {
    std::string maVach;           // 8 hoac 13 chu so, khoa B-Tree
    std::string ten, nhom;  double donGia;
    int tonKho;                   // = tong soLuong cac lo
    int nguongCanhBao;  LotQueue dsLo;  int demTam;   // demTam: bien tam cho bao cao
};

struct BTreeNode {
    int n; bool leaf;
    MatHang** keys;               // toi da m-1 khoa (+1 o de tach)
    BTreeNode** child;            // toi da m con (+1)
};
struct BTree { BTreeNode* root; int m; int size;
               long long soNutTruyXuat; bool inKhiTachGop; };
\end{lstlisting}

\section{Lựa chọn cấu trúc dữ liệu}

\subsection{Chỉ mục mặt hàng: B-Tree so với AVL, bảng băm và mảng sắp xếp}

Thao tác thường xuyên nhất là tìm mặt hàng theo mã vạch (mỗi lần nhập, xuất, sửa, báo cáo top bán chạy).
Ngoài ra cần chèn, xóa và duyệt toàn bộ theo thứ tự mã vạch (để ghi tệp và lập báo cáo).
Bảng~\ref{tab:sosanhchimuc} so sánh các phương án với $n$ mặt hàng.

\begin{table}[htbp]
\centering\small
\caption{So sánh các cấu trúc làm chỉ mục mặt hàng ($n$ khóa)}
\label{tab:sosanhchimuc}
\begin{tabularx}{\linewidth}{@{}l c c c Y@{}}
\toprule
\textbf{Cấu trúc} & \textbf{Tìm} & \textbf{Chèn / xóa} & \textbf{Duyệt có thứ tự} & \textbf{Nhận xét} \\
\midrule
Mảng sắp xếp & $O(\log n)$ & $O(n)$ & $O(n)$ & Dời phần tử khi chèn/xóa, chậm khi kho thay đổi nhiều \\
Bảng băm & $O(1)$ TB & $O(1)$ TB & $O(n\log n)$ & Không có thứ tự; phải sắp lại khi ghi tệp/báo cáo; xấu nhất $O(n)$ \\
Cây AVL & $O(\log_2 n)$ & $O(\log_2 n)$ & $O(n)$ & Mỗi nút 1 khóa, cao $\approx 1{,}44\log_2 n$ trong trường hợp xấu \\
\textbf{B-Tree bậc $m$} & $O(\log_{\lceil m/2\rceil} n)$ nút & $O(m\log_{\lceil m/2\rceil} n)$ & $O(n)$ & Mỗi nút tới $m-1$ khóa, cây thấp, ít nút phải truy xuất \\
\bottomrule
\end{tabularx}
\end{table}

\paragraph{So sánh trực tiếp B-Tree với AVL.} Cả hai đều là cây tìm kiếm cân bằng với chi phí $O(\log n)$,
nhưng khác nhau ở \emph{số nút phải đi qua}. Với $n = 10^5$:
\begin{itemize}
  \item AVL có chiều cao ít nhất $\lceil\log_2(n+1)\rceil = 17$ và nhiều nhất
        $1{,}44\log_2(n+2) - 0{,}33 \approx 23{,}6$ mức~\cite{avl,knuth}; mỗi lần tìm đi qua trung bình xấp xỉ
        $\log_2 n \approx 16{,}6$ nút.
  \item B-Tree bậc $m = 10$ có chiều cao không quá $1 + \log_5\frac{n+1}{2} \approx 7{,}7$; thực nghiệm ở
        Chương~\ref{chap:thucnghiem} đo được chiều cao 6 và trung bình \textbf{5,84 nút} mỗi lần tìm.
\end{itemize}
Đổi lại, mỗi nút B-Tree chứa nhiều khóa nên phải tìm nhị phân trong nút ($O(\log m)$ phép so sánh), và
thao tác chèn/xóa phức tạp hơn (tách, mượn, gộp nút thay vì xoay cây). Khi toàn bộ dữ liệu nằm trong bộ nhớ,
hai cấu trúc có tốc độ cùng bậc; ưu điểm thật sự của B-Tree xuất hiện khi mỗi lần truy xuất nút là một lần đọc
khối đĩa -- trường hợp của kho dữ liệu thực tế với hàng triệu mặt hàng~\cite{bayer,comer}. B-Tree cũng có lợi
về bộ nhớ đệm (cache) vì các khóa trong một nút nằm liền nhau. Đây là lý do đề tài chọn B-Tree, đồng thời
cho phép thực nghiệm ảnh hưởng của bậc $m$.

\subsection{Danh sách lô của mặt hàng: hàng đợi DSLK so với mảng vòng}

Các lô của một mặt hàng được xử lý đúng theo thứ tự FIFO: nhập kho thêm lô vào \emph{cuối}, xuất kho lấy từ
\emph{đầu}. Đây chính là hàng đợi. Hai cách cài đặt được cân nhắc:

\begin{table}[htbp]
\centering\small
\caption{So sánh cách cài đặt hàng đợi lô}
\label{tab:sosanhhangdoi}
\begin{tabularx}{\linewidth}{@{}l Y Y@{}}
\toprule
& \textbf{DSLK đơn có \code{front}, \code{rear}} (được chọn) & \textbf{Mảng vòng} \\
\midrule
Thêm cuối / lấy đầu & $O(1)$ & $O(1)$, nhưng $O(k)$ khi đầy phải cấp phát lại và chép \\
Số lô tối đa & Không giới hạn & Cố định theo dung lượng hoặc phải mở rộng \\
Bộ nhớ & Mỗi lô thêm 1 con trỏ \code{next} & Liền khối, có ô trống dự phòng \\
Duyệt các lô (báo cáo HSD) & $O(k)$ & $O(k)$, thân thiện cache hơn \\
Hàng nghìn mặt hàng ít lô & Chỉ cấp phát đúng số lô thật & Mỗi mặt hàng giữ một mảng dự phòng \\
\bottomrule
\end{tabularx}
\end{table}

Số lô mỗi mặt hàng thay đổi liên tục và thường nhỏ, trong khi số mặt hàng lớn. DSLK không lãng phí bộ nhớ dự
phòng cho từng mặt hàng, không cần cấp phát lại, và việc gỡ lô đã xuất hết chỉ là dời con trỏ \code{front}.

\subsection{Các cấu trúc phụ}

\begin{itemize}
  \item \textbf{Danh sách phiếu}: DSLK đơn có con trỏ \code{tail} để thêm phiếu mới vào cuối trong $O(1)$;
        phiếu được duyệt tuần tự khi in danh sách, lọc theo ngày và tính top bán chạy.
  \item \textbf{Mảng động \code{DsMatHang}} (thay \code{std::vector}): dùng cho kết quả tìm kiếm và báo cáo, vì
        Merge Sort cần truy cập ngẫu nhiên. Mảng gấp đôi dung lượng khi đầy nên thêm phần tử có chi phí khấu
        trừ $O(1)$.
  \item \textbf{Merge Sort} (thay \code{std::sort}): được chọn thay vì Quick Sort vì \emph{ổn định} -- cần cho
        sắp xếp nhiều tiêu chí như ``theo nhóm, cùng nhóm thì theo tên'' -- và luôn $O(n\log n)$, không có trường hợp
        xấu $O(n^2)$. Đổi lại cần $O(n)$ bộ nhớ phụ.
\end{itemize}

\section{Sơ đồ mô-đun}

Chương trình gồm 12 mô-đun (Hình~\ref{fig:modun}). Mũi tên $A \to B$ nghĩa là $A$ gọi hàm của $B$.
Các mô-đun đều dùng kiểu dữ liệu trong \code{Types.h} và hàm ngày trong \code{Date}.

\begin{figure}[htbp]
\centering
\resizebox{\linewidth}{!}{%
\begin{tikzpicture}[
  mod/.style={draw,rounded corners=2pt,minimum width=2.25cm,minimum height=0.95cm,align=center,font=\small},
  dep/.style={-{Stealth[length=2mm]},thick,black!60}]
  \node[mod,fill=cTin,minimum width=3cm] (main) at (6.5,0) {\textbf{main.cpp}\\[-2pt]{\scriptsize menu}};
  \node[mod,fill=cTin]    (val) at (0,-2.2)    {\textbf{Validate}\\[-2pt]{\scriptsize kiểm tra nhập}};
  \node[mod,fill=cTin]    (fio) at (2.6,-2.2)  {\textbf{FileIO}\\[-2pt]{\scriptsize đọc/ghi tệp}};
  \node[mod,fill=cTan]    (inv) at (5.2,-2.2)  {\textbf{Inventory}\\[-2pt]{\scriptsize nhập/xuất kho}};
  \node[mod,fill=cVu]     (rep) at (7.8,-2.2)  {\textbf{Report}\\[-2pt]{\scriptsize tìm, báo cáo}};
  \node[mod,fill=cTin]    (gen) at (10.4,-2.2) {\textbf{DataGen}\\[-2pt]{\scriptsize sinh dữ liệu}};
  \node[mod,fill=cNguyen] (ben) at (13,-2.2)   {\textbf{Benchmark}\\[-2pt]{\scriptsize thực nghiệm}};
  \node[mod,fill=cTan]    (lq)  at (3.4,-4.4)  {\textbf{LotQueue}\\[-2pt]{\scriptsize hàng đợi lô}};
  \node[mod,fill=cTan]    (ph)  at (6,-4.4)    {\textbf{Phieu}\\[-2pt]{\scriptsize DS phiếu}};
  \node[mod,fill=cNguyen] (bt)  at (8.6,-4.4)  {\textbf{BTree}\\[-2pt]{\scriptsize chỉ mục mã vạch}};
  \node[mod,fill=cVu]     (so)  at (11.2,-4.4) {\textbf{Sort}\\[-2pt]{\scriptsize Merge Sort}};
  \node[mod,fill=cChung,minimum width=7cm] (ty) at (3.8,-6.2) {\textbf{Types.h}\\[-2pt]{\scriptsize struct dùng chung}};
  \node[mod,fill=cTin,minimum width=7cm]   (da) at (11.2,-6.2) {\textbf{Date}\\[-2pt]{\scriptsize ngày: kiểm tra, cộng/trừ, so sánh}};
  \foreach \t in {val,fio,inv,rep,gen,ben} \draw[dep] (main.south) -- (\t.north);
  \draw[dep] (fio) -- (val);
  \draw[dep] (fio) -- (inv);
  \draw[dep] (rep) -- (inv);
  \draw[dep] (inv.south) -- (lq.north);
  \draw[dep] (inv.south) -- (ph.north);
  \draw[dep] (inv.south) -- (bt.north);
  \draw[dep] (rep.south) -- (bt.north);
  \draw[dep] (rep.south) -- (so.north);
  \draw[dep] (ben.south) -- (bt.north);
  \draw[dashed,thick,black!45] (-1.3,-5.4) -- (14.4,-5.4);
  \node[font=\scriptsize\itshape,anchor=south west] at (-1.3,-5.4) {mọi mô-đun phía trên dùng};
  % chu thich mau
  \node[font=\scriptsize,anchor=west] at (-1.2,-7.4) {Phụ trách:};
  \node[draw,fill=cNguyen,minimum size=3mm] at (0.6,-7.4) {}; \node[font=\scriptsize,anchor=west] at (0.8,-7.4) {\TVa};
  \node[draw,fill=cTan,minimum size=3mm] at (3.2,-7.4) {};    \node[font=\scriptsize,anchor=west] at (3.4,-7.4) {\TVb};
  \node[draw,fill=cVu,minimum size=3mm] at (5.6,-7.4) {};     \node[font=\scriptsize,anchor=west] at (5.8,-7.4) {\TVc};
  \node[draw,fill=cTin,minimum size=3mm] at (7.9,-7.4) {};    \node[font=\scriptsize,anchor=west] at (8.1,-7.4) {\TVd};
  \node[draw,fill=cChung,minimum size=3mm] at (10.3,-7.4) {}; \node[font=\scriptsize,anchor=west] at (10.5,-7.4) {Cả nhóm};
\end{tikzpicture}}
\caption{Sơ đồ mô-đun và quan hệ gọi hàm}
\label{fig:modun}
\end{figure}

\section{Tổ chức dữ liệu trong bộ nhớ}

Hình~\ref{fig:bonho} minh họa trạng thái bộ nhớ khi kho chứa mặt hàng \code{8930000000002} (Mì gói Loại 2) với
hai lô còn hàng. Khóa trong nút B-Tree là con trỏ tới bản ghi \code{MatHang}; bản ghi giữ hàng đợi lô với
\code{front} là lô cũ nhất. Danh sách phiếu độc lập với cây và chỉ tham chiếu mặt hàng qua mã vạch.

\begin{figure}[htbp]
\centering
\resizebox{\linewidth}{!}{%
\begin{tikzpicture}[
  bnode/.style={rectangle split,rectangle split horizontal,rectangle split parts=3,draw,font=\scriptsize\ttfamily,
                inner sep=2.5pt,fill=cNguyen},
  rec/.style={draw,align=left,font=\scriptsize\ttfamily,fill=cTan,inner sep=4pt},
  lo/.style={rectangle split,rectangle split horizontal,rectangle split parts=4,draw,font=\scriptsize\ttfamily,
             fill=cTan!60,inner sep=2.5pt},
  ph/.style={rectangle split,rectangle split horizontal,rectangle split parts=2,draw,font=\scriptsize\ttfamily,
             fill=cTan!30,inner sep=2.5pt},
  ptr/.style={-{Stealth[length=2mm]},thick}]
  \node[bnode] (r) at (0,0) {\nodepart{one}...009\nodepart{two}...018\nodepart{three}...027};
  \node[bnode] (c1) at (-2.4,-1.4) {\nodepart{one}...003\nodepart{two}...006\nodepart{three}\phantom{000}};
  \node[bnode] (l1) at (-4.6,-2.8) {\nodepart{one}...001\nodepart{two}...002\nodepart{three}\phantom{000}};
  \draw[ptr] ($(r.south west)!0.1!(r.south east)$) -- (c1.north);
  \draw[ptr] ($(c1.south west)!0.1!(c1.south east)$) -- (l1.north);
  \node[font=\scriptsize\itshape,anchor=west] at (1.9,0) {B-Tree $m=5$ (35 khóa, cao 3)};
  \node[rec] (mh) at (-1.2,-4.6) {MatHang\\maVach = 8930000000002\\ten = "Mi goi Loai 2"\\tonKho = 64 \ \ (= 25 + 39)\\dsLo: front, rear, count = 2};
  \draw[ptr] (l1.two south) to[out=-90,in=180] (mh.west);
  \node[lo] (lo4) at (4.2,-4.1) {\nodepart{one}LO4\nodepart{two}09/11/2026\nodepart{three}HSD 15/04/2027\nodepart{four}25};
  \node[lo] (lo5) at (4.2,-5.2) {\nodepart{one}LO5\nodepart{two}16/11/2026\nodepart{three}HSD 18/04/2027\nodepart{four}39};
  \node[font=\scriptsize,anchor=west] (nul) at (7.0,-5.2) {nullptr};
  \draw[ptr] (mh.east |- lo4.west) -- node[above,font=\scriptsize]{front} (lo4.west);
  \draw[ptr] (mh.east |- lo5.west) -- node[below,font=\scriptsize]{rear} (lo5.west);
  \draw[ptr] ([xshift=6mm]lo4.south west) -- node[right,font=\scriptsize]{next} ([xshift=6mm]lo5.north west);
  \draw[ptr] (lo5.east) -- (nul.west);
  \node[font=\scriptsize\itshape] at (4.2,-3.5) {hàng đợi lô (cũ $\to$ mới)};
  % danh sach phieu
  \node[ph] (p1) at (-3.6,-7) {\nodepart{one}PN00001\nodepart{two}...001};
  \node[ph] (p2) at (-1.2,-7) {\nodepart{one}PN00002\nodepart{two}...001};
  \node[font=\small] (pd) at (0.6,-7) {$\cdots$};
  \node[ph] (p3) at (2.4,-7) {\nodepart{one}PX00099\nodepart{two}...035};
  \node[font=\scriptsize] (n2) at (4.4,-7) {nullptr};
  \draw[ptr] (p1) -- (p2); \draw[ptr] (p2) -- (pd); \draw[ptr] (pd) -- (p3); \draw[ptr] (p3) -- (n2);
  \node[font=\scriptsize] (hd) at (-3.6,-8) {head};  \draw[ptr] (hd) -- (p1);
  \node[font=\scriptsize] (tl) at (2.4,-8) {tail};   \draw[ptr] (tl) -- (p3);
  \node[font=\scriptsize\itshape,anchor=west] at (5.4,-7) {DanhSachPhieu (99 phiếu)};
\end{tikzpicture}}
\caption{Tổ chức dữ liệu trong bộ nhớ: B-Tree, bản ghi mặt hàng, hàng đợi lô và danh sách phiếu}
\label{fig:bonho}
\end{figure}

\section{Định dạng tệp và mã kết quả}

Dữ liệu lưu trong ba tệp văn bản, mỗi dòng một bản ghi, các trường phân cách bằng ký tự \code{|}, ngày dạng
\code{dd/mm/yyyy} (Bảng~\ref{tab:tep}). Trong tệp, mã lô và mã phiếu có tiền tố cho dễ đọc; trong chương trình
chúng là số nguyên: \code{LO00001}~$\leftrightarrow$~1, \code{PN00006} (phiếu nhập) và \code{PX00006}
(phiếu xuất)~$\leftrightarrow$~6. Phiếu nhập và xuất dùng chung một dãy số, nên không có hai phiếu cùng số.
Quy đổi chỉ được định nghĩa ở một nơi (\code{tachMaLo}, \code{taoMaLo}, \code{tachMaPhieu}, \code{taoMaPhieu}
trong \code{FileIO}).

\begin{table}[htbp]
\centering\small
\caption{Định dạng ba tệp dữ liệu}
\label{tab:tep}
\begin{tabularx}{\linewidth}{@{}l Y@{}}
\toprule
\textbf{Tệp} & \textbf{Định dạng dòng và ví dụ} \\
\midrule
\code{mathang.txt} & \code{maVach|ten|nhom|donGia|nguongCanhBao} \\
 & \code{8930000000002|Mi goi Loai 2|Thuc pham kho|110000|10} \\
\code{lo.txt} & \code{maVach|maLo|ngayNhap|hsd|soLuong|giaNhap} \quad (số lượng \emph{còn lại}, theo thứ tự FIFO) \\
 & \code{8930000000002|LO00004|09/11/2026|15/04/2027|25|82500} \\
\code{phieu.txt} & \code{maPhieu|loai|ngay|maVach|soLuong|donGia} \\
 & \code{PX00006|XUAT|16/11/2026|8930000000002|51|110000} \\
\bottomrule
\end{tabularx}
\end{table}

Các hàm nghiệp vụ trả về một giá trị \code{enum KetQua} thay vì in thông báo trực tiếp, để menu quyết định cách
hiển thị và để kiểm thử tự động được: \code{KQ\_OK}, \code{KQ\_TRUNG\_MA}, \code{KQ\_KHONG\_TIM\_THAY},
\code{KQ\_KHONG\_DU\_HANG}, \code{KQ\_DU\_LIEU\_SAI}, \code{KQ\_CON\_TON\_KHO}, \code{KQ\_LOI\_TEP}.
